% !TeX program = xelatex
% !TeX encoding = UTF-8
\documentclass[12pt]{beamer}
\IfFileExists{template_style.tex}{% Style adapted from the user-provided output/presentations/main.tex.
\usetheme{Madrid}
\usefonttheme{serif}

\usepackage[english]{babel}
\usepackage{amsmath,amssymb,mathtools}
\usepackage{tikz}
\usetikzlibrary{arrows.meta,positioning,calc}

\definecolor{accent}{RGB}{55,55,180}
\definecolor{diagramblue}{RGB}{48,103,166}
\definecolor{diagramred}{RGB}{178,65,56}
\definecolor{diagramgreen}{RGB}{47,128,92}
\definecolor{softgray}{RGB}{118,118,118}
\definecolor{lightgray}{RGB}{224,224,224}

\setbeamercolor{structure}{fg=accent}
\setbeamercolor{frametitle}{fg=white,bg=accent}
\setbeamercolor{title}{fg=white,bg=accent}
\setbeamercolor{palette primary}{fg=white,bg=accent}
\setbeamercolor{palette secondary}{fg=white,bg=accent!88!black}
\setbeamercolor{palette tertiary}{fg=white,bg=accent!74!black}
\setbeamercolor{palette quaternary}{fg=white,bg=accent!62!black}
\setbeamercolor{normal text}{fg=black,bg=white}
\setbeamertemplate{navigation symbols}{}
\setbeamersize{text margin left=0.68cm,text margin right=0.68cm}

\newcommand{\eyebrow}[1]{%
  {\color{accent}\bfseries #1}\par
  \vspace{0.04cm}{\color{lightgray}\rule{\linewidth}{0.45pt}}\vspace{0.10cm}}
\newcommand{\keyterm}[1]{\textcolor{accent}{\bfseries #1}}
\newcommand{\dd}{\,\mathrm d}
}{% Style adapted from the user-provided output/presentations/main.tex.
\usetheme{Madrid}
\usefonttheme{serif}

\usepackage[english]{babel}
\usepackage{amsmath,amssymb,mathtools}
\usepackage{tikz}
\usetikzlibrary{arrows.meta,positioning,calc}

\definecolor{accent}{RGB}{55,55,180}
\definecolor{diagramblue}{RGB}{48,103,166}
\definecolor{diagramred}{RGB}{178,65,56}
\definecolor{diagramgreen}{RGB}{47,128,92}
\definecolor{softgray}{RGB}{118,118,118}
\definecolor{lightgray}{RGB}{224,224,224}

\setbeamercolor{structure}{fg=accent}
\setbeamercolor{frametitle}{fg=white,bg=accent}
\setbeamercolor{title}{fg=white,bg=accent}
\setbeamercolor{palette primary}{fg=white,bg=accent}
\setbeamercolor{palette secondary}{fg=white,bg=accent!88!black}
\setbeamercolor{palette tertiary}{fg=white,bg=accent!74!black}
\setbeamercolor{palette quaternary}{fg=white,bg=accent!62!black}
\setbeamercolor{normal text}{fg=black,bg=white}
\setbeamertemplate{navigation symbols}{}
\setbeamersize{text margin left=0.68cm,text margin right=0.68cm}

\newcommand{\eyebrow}[1]{%
  {\color{accent}\bfseries #1}\par
  \vspace{0.04cm}{\color{lightgray}\rule{\linewidth}{0.45pt}}\vspace{0.10cm}}
\newcommand{\keyterm}[1]{\textcolor{accent}{\bfseries #1}}
\newcommand{\dd}{\,\mathrm d}
}

% Five content slides, using the supplied Madrid theme and 4:3 canvas.
\setbeamerfont{frametitle}{size=\fontsize{17}{20}\selectfont}
\setlength{\abovedisplayskip}{6pt}
\setlength{\belowdisplayskip}{6pt}
\setlength{\abovedisplayshortskip}{3pt}
\setlength{\belowdisplayshortskip}{4pt}
\setbeamertemplate{itemize items}[circle]
\newcommand{\R}{\mathbb R}
\newcommand{\End}{\operatorname{End}}
\newcommand{\Pp}{\mathcal P}
\newcommand{\sourcecredit}[2]{%
  {\fontsize{8}{10}\selectfont\color{softgray}%
  Student Assessment Exercise, Project #1, item #2 (p.~12).}\par\smallskip}

\title[Algebra and Linear Dynamics]{Algebraic Structures and Linear Dynamics}
\author[H. Zhang]{Heze Zhang}
\date{November 6, 2026}

\begin{document}

\section{Algebra and spaces}

\begin{frame}[t]{State space and observable algebras}
  \keyterm{Measurable state space}
  \begin{equation*}
    V=\R^d,\qquad \mathcal B(V)=\sigma(\text{Euclidean open sets}).
  \end{equation*}
  \vspace{-2mm}
  \begin{columns}[T,onlytextwidth]
    \begin{column}{0.48\textwidth}
      \keyterm{Linear maps on states}
      \begin{equation*}
        \End(V)\cong M_d(\R).
      \end{equation*}
      Addition and scalar\\multiplication are pointwise.
      \begin{equation*}
        AB=A\circ B.
      \end{equation*}
      Associative and unital, with identity $I_V$.
    \end{column}
    \begin{column}{0.48\textwidth}
      \keyterm{Functions on states}
      \begin{equation*}
        \mathcal A=B_b(V).
      \end{equation*}
      Bounded Borel measurable real functions, with
      \begin{equation*}
        (fg)(x)=f(x)g(x).
      \end{equation*}
      Commutative and unital, with identity $\mathbf1$.
    \end{column}
  \end{columns}
  \note{Foundation: paper/main.tex, Section 1.1. Linear transformations form an algebra
  under composition. Observable functions form an algebra under pointwise multiplication.
  End(V) is generally noncommutative when d is at least two. Physical observables
  may be unbounded and are then used under appropriate integrability assumptions.}
\end{frame}

\begin{frame}[t]{Observable and probability evolution}
  \keyterm{Probability laws and the integration pairing}
  \begin{equation*}
    \mu\in\Pp(V),\qquad
    \langle f,\mu\rangle=\int_V f\dd\mu=\mathbb E[f(X)],\quad X\sim\mu.
  \end{equation*}
  For a Borel measurable transformation $F:V\to V$, define
  \begin{equation*}
    \begin{aligned}
      \mathsf U_F f&=f\circ F,\\
      \mathsf T_F\mu&=F_\#\mu,\quad
        (F_\#\mu)(C)=\mu(F^{-1}(C)),\quad C\in\mathcal B(V).
    \end{aligned}
  \end{equation*}
  \begin{equation*}
    \boxed{\langle\mathsf U_Ff,\mu\rangle
         =\langle f,\mathsf T_F\mu\rangle.}
  \end{equation*}
  $\Pp(V)$ is convex. The pairing extends bilinearly to finite signed measures.
  \note{Foundation: the manuscript's observable--measure duality. The pairing applies
  to bounded Borel f and requires no density. U_F acts on observables, T_F on measures.
  A nonlinear state map still induces a linear operator on observables. The matrix
  U(t,s) used later is instead the homogeneous evolution on the state space.}
\end{frame}

\section{Project 9.1: linear dynamics}

\begin{frame}[t]{Project 9.1: the ODE and its solution spaces}
  \sourcecredit{9}{1}
  Let $A:I\to\End(V)$ and $B:I\to V$ be continuous. Solve
  \begin{equation*}
    \dot x_t=A(t)x_t+B(t),\qquad x_s=x.
  \end{equation*}
  The ODE theorem gives a unique trajectory on $I$.
  \medskip

  \keyterm{A linear operator on trajectories}
  \begin{equation*}
    \begin{aligned}
      \mathcal X&=C^1(I,V),\qquad \mathcal Y=C(I,V),\\
      \mathcal L:\mathcal X\to\mathcal Y,\qquad
      (\mathcal L\gamma)(t)&=\dot\gamma(t)-A(t)\gamma(t).
    \end{aligned}
  \end{equation*}
  The homogeneous solutions form the vector space
  \begin{equation*}
    \mathcal H=\ker\mathcal L
      =\{\eta\in\mathcal X:\dot\eta=A\eta\}.
  \end{equation*}
  \note{Source: docs/StudentAssessmentExercise_unlocked.pdf, Project 9, item 1, p.12.
  On each compact interval J containing s, continuity bounds A and B and gives a
  uniform global Lipschitz constant in the state and a linear-growth bound.
  The ODE theorem gives solutions on J, and uniqueness joins them on overlaps.
  Hence the entire given time interval I is covered. This applies the earlier
  existence theorem rather than repeating its proof.}
\end{frame}

\begin{frame}[t]{Homogeneous evolution}
  Evaluation at time $s$ identifies trajectories with states:
  \begin{equation*}
    E_s:\mathcal H\to V,\quad \eta\mapsto\eta(s),\qquad
    \mathcal H\cong V,\quad\dim\mathcal H=d.
  \end{equation*}
  Existence makes $E_s$ onto. Uniqueness makes it one-to-one.
  \medskip

  \keyterm{Evolution on states}
  \begin{equation*}
    \begin{aligned}
      U(t,s)&=E_tE_s^{-1},\qquad U(s,s)=I_V,\\
      \partial_tU(t,s)&=A(t)U(t,s),\\
      U(t,r)U(r,s)&=U(t,s),\qquad U(t,s)^{-1}=U(s,t).
    \end{aligned}
  \end{equation*}
  \note{E_s is a linear isomorphism by existence and uniqueness of homogeneous
  solutions. Its inverse recovers a complete trajectory from the initial state.
  Composition and inverses follow by cancellation of these evaluation maps.
  For constant A, U(t,s)=exp((t-s)A). For noncommuting time-dependent A, the
  fundamental matrix is time ordered. No ordinary exponential is assumed here.}
\end{frame}

\begin{frame}[t]{Variation of constants and verification}
  Seek a particular solution $\gamma_p(t)=U(t,s)c(t)$ with $c(s)=0$.
  \begin{equation*}
    \begin{aligned}
      \dot\gamma_p&=A(t)\gamma_p+U(t,s)\dot c,\\
      U(t,s)\dot c(t)=B(t)
        &\ \Longrightarrow\ \dot c(t)=U(s,t)B(t).
    \end{aligned}
  \end{equation*}
  Integrate and use composition to obtain
  \begin{equation*}
    \boxed{x_t=U(t,s)x+\int_s^t U(t,r)B(r)\dd r.}
  \end{equation*}
  \keyterm{Verification}\quad $U(s,s)=I_V$ gives $x_s=x$. The Leibniz rule gives
  \begin{equation*}
    \dot x_t=A(t)\left[U(t,s)x+\int_s^t U(t,r)B(r)\dd r\right]+B(t)
            =A(t)x_t+B(t).
  \end{equation*}
  \note{Integrating c prime from s to t gives c(t)=integral U(s,r)B(r)dr.
  Multiplication by U(t,s) and the composition identity give the particular
  solution gamma_p(t)=integral U(t,r)B(r)dr. The full solution set with fixed B
  is the affine space gamma_p+H. The initial condition selects one trajectory.
  Differentiation under the integral is valid by joint continuity of U and its
  terminal-time derivative on compact intervals. Cooling and the detailed
  time-ordering discussion remain in the manuscript rather than these five slides.}
\end{frame}

\end{document}
